\begin{comment}


\section{First Step}

IDEE DE PLAN/TRAJECTOIRE (à partir de la section 6 pour avoir la vision):
\begin{itemize}
    \item 1- \textbf{Comment on a crée/ quelles sont nos méthode EWC :} On pose notre Loss EWC (inspiré du paper) -> methodes découlant de cette loss: méthode A et B (bien appuyer sur sequentialité step 1 puis step 2) 
    
    \item 2- \textbf{Step 1:} overlap $m_1$, transition $\lambda_1 >1$, formulation théorique de notre estimateur (garder uniquement $\hat x^{(1)}_{P_1}$ les autres en annexe peut etre)
    
    \item 3- \textbf{En préparation du Step 2:} Etude de la matrice Hessienne A évalué en $\hat x^{(1)}_{P_1}$: 
    
    Eigenvalues, transition, overlaps avec signaux
    
    (motivation: elle est centrale dans nos deux méthodes étant donné que c'est la Hessienne de notre Loss)
    \textit{Ici on ne parle pas encore d'estimateur du step 2, uniquement d'eigenvalue et eigenvector}

    \item 4- \textbf{Methode A:}
    Estimateur final au Step 2: 
    
    top(A), condition pour retrouver le signal ($\theta_{max}>1$), overlapp avec les deux signaux

    \item 5- \textbf{Methdode B:}
    Estimateur final: ecrire l'expression dans la base de A, condition pour retrouver le signal ($\theta_{max}>1$ au moins 1 spike qui sort du bulk -> au moins un overlap non nul), overlapp avec les deux signaux

\end{itemize}

Recover the first signal $\vect{x_1}$ Let us define our fist-step estimator $\hat{\vect x}^{(1)}$ as the top eigenvector of the first observaation $\rdmmat Y_1 = \lambda_1 \vect x_1 \vect x_1 ^T + \rdmmat Z_1$ 
\subsection{Ideal case: complete recovery}
We consider that in the first step our estimator of the signal completely recovers it.
For the approximation of a perfect recovery of the first signal ($ \hat{\vect x}^{(1)} = \vect x_1$), the Hessian is of the form $\rdmmat A = \rdmmat P + \rdmmat W + c \rdmmat I_n$ where:
\[
\begin{cases}
    \rdmmat P = a \vect x_1 \vect x_1^T + b \vect x_2 \vect x_2^T \qquad \text{with} \quad a = 2\mu_R \Big((\lambda_1-1)^2 + \frac{1}{n} \Big) > 0, \quad b = \lambda_2 \\
    \rdmmat W = \rdmmat Z _2\\
    c = \Big(\frac{2\mu_R}{n} - 1 \Big) 
\end{cases}
\]
The recovery transitions are then given by the general results of section \ref{sec:general_results}.

\subsection{Non-Ideal case: partial recovery}
We consider a way to approximate the first step estimator $\hat{\vect x}^{(1)}$.

\subsubsection{Estimation via first time step overlap}
The first time step overlap with the signal $m_1$, represents the recovered fraction of the signal $\vect x_1$. An possible approximation of our estimator is then:

\begin{equation}
\boxed{\;
    \hat{\vect x}^{(1)}_{m_1} = m_1 \vect x_1\;
    }
    \label{eq:x_hat_m1}
\end{equation}
Then as for the perfect recovery, the Hessian is of the form $\rdmmat A = \rdmmat P + \rdmmat W + c \rdmmat I_n$.
The recovery transitions are then given by Eq. \ref{eq:BBP_overlap_k_spikes} and Eq. \ref{eq:BBP-k-spikes}, with $k=2$ and the parameters:
\[
\begin{cases}
     a = 2\mu_R \Big((\lambda_1-1)^2 + \frac{1}{n} \Big) m_1^2  \\
     b = \lambda_2 \\
    c = \Big(\frac{2\mu_R}{n} - 1 \Big) 
\end{cases}
\]

\subsubsection{Estimation with an orthogonal projector of the signal}
A more complete way to approximate our estimator is to also consider the not recovered part of our signal. Considering a rank one spike $\vect g \in \R^n$ with $g_i \sim \mathcal{N}( 0, 1)$, project on the orthogonal subspace of the signal $\vect x_1$, the estimator can be approximate by:
\begin{equation}
\boxed{
    \hat{\; \vect x}^{(1)}_{P_1} = m_1 \vect x_1 + \sqrt{1 - m_1 ^2} \frac{\rdmmat P_{\perp x_1}\vect g}{\norm{P_{\perp x_1}\vect g}} \;} \quad \text{with} \quad \rdmmat P_{\perp x_1} = \rdmmat I - \vect x_1 \vect x_1^\top
    \label{eq:x_hat_p1}
\end{equation}
defining a rotated spike $\vect x_3 = \frac{\rdmmat P_{\perp x_1}\vect g}{\norm{P_{\perp x_1}\vect g}} $, the Hessian matrix $\rdmmat A$ is given, as before, by $ \rdmmat A = \rdmmat P + \rdmmat W + c \rdmmat I_n$ with:
\[
\rdmmat P = \alpha_{11} \vect x_1 \vect x_1^T + \alpha_{22} \vect x_2 \vect x_2^T + \alpha_{33} \vect x_3 \vect x_3^T + \alpha_{13}\big( \vect x_1 \vect x_3^T + \vect x_3 \vect x_1^T \big)
\]
and the coefficients:
\[
\begin{cases}
     \alpha_{11} = 2\mu_R \Big((\lambda_1-1)^2 + \frac{1}{n} \Big) m_1^2  \\
     \alpha_{22} = \lambda_2 \\
     \alpha_{33} = 2\mu_R \Big((\lambda_1-1)^2 + \frac{1}{n} \Big)(1 -m_1^2) \\
     \alpha_{13} = 2\mu_R \Big((\lambda_1-1)^2 + \frac{1}{n} \Big) m_1 \sqrt{1- m_1^2}
\end{cases}
\]
The orthonormal basis of $(\vect e_1, \vect e_2, \vect e_3)$ of $\rdmmat P$ is:
\[
\begin{cases}
    \vect x_1 = \vect e_1 \\
    \vect x_2 = \rho \vect e_1 + \Tilde{\rho} \vect e_3 + \sqrt{1 - \rho^2 - \Tilde{\rho}^2} \vect e_2 \\
    \vect x_3 = \vect e_3
\end{cases}
\]
Where $\rho$ is the imposed correlation between signal 1 and signal 2 ($\lvert \langle \vect x_1, \vect x_2 \rangle \rvert \xrightarrow[n\to\infty]{\text{a.s.}} \rho $) and $\Tilde{\rho}$ is the induced correlation between the second signal and our created third spike ($\lvert \langle \vect x_2, \vect x_3 \rangle \rvert = \Tilde{\rho} \rangle \rvert \approx \frac{1}{\sqrt n} \xrightarrow[n\to\infty]{\text{a.s.}} 0$).
In this basis, the matrix $\rdmmat P$ is then:
\begin{comment}
\[
\begin{aligned}
\rdmmat{P} = &{}
(\alpha_{11}+\alpha_{22}\rho^2)\,
\vect e_1 \vect e_1^T
+\alpha_{22}(1-\rho^2-\widetilde{\rho}^{\,2})\,
\vect e_2 \vect e_2^T
 +(\alpha_{33}+\alpha_{22}\widetilde{\rho}^{\,2})\,
\vect e_3 \vect e_3^T
\\
&+\alpha_{22}\rho\sqrt{1-\rho^2-\widetilde{\rho}^{\,2}}\,
\left(\vect e_1 \vect e_2^T+\vect e_2 \vect e_1^T\right)
+(\alpha_{13}+\alpha_{22}\rho\widetilde{\rho})\,
\left(\vect e_1 \vect e_3^T+\vect e_3 \vect e_1^T\right)
+\alpha_{22}\widetilde{\rho}\sqrt{1-\rho^2-\widetilde{\rho}^{\,2}}\,
\left(\vect e_2 \vect e_3^T+\vect e_3 \vect e_2^T\right). 
\end{aligned}
\]
Leading to:

\[ 
\rdmmat P =  \sum_{i,j = 1}^3 \beta_{ij} \vect e_i \vect e_j^T
\]
where the coefficients $\beta_{ij} \equiv \beta_{ij}(\lambda_1, m_1, \lambda_2, \mu_R,\rho, \Tilde{\rho})$ are given in the annex table \ref{tab:coefficients_beta_A(xP1)}.
\paragraph{Remark} As the third spike $\vect x_3$ has no preferred direction except it's orthogonality with $\vect x_1$, for $n$ large enough, it's overlap with the second signal $\vect x_2$ tend to zero.
\[
    \Tilde{\rho}= \lvert \langle \vect x_2, \vect x_3 \rangle \rvert \approx \frac{1}{\sqrt n} \xrightarrow[n\to\infty]{\text{a.s.}} 0
\]
In the limit $\widetilde{\rho}\to 0$, the matrix
Let \(\rdmmat B\) be the real symmetric matrix
\[
\rdmmat B=
\begin{pmatrix}
\beta_{11}&\beta_{12}&\beta_{13}\\
\beta_{12}&\beta_{22}& 0\\
\beta_{13}& 0 &\beta_{33}
\end{pmatrix}.
\]
has entries detailed in T\ref{truc}. And it's top eigenvalue is:

\[
\begin{aligned}
\beta_{11}&=\alpha_{11}+\alpha_{22}\rho^2, &
\beta_{22}&=\alpha_{22}(1-\rho^2), &
\beta_{33}&=\alpha_{33},\\
\beta_{12}&=\alpha_{22}\,\rho\sqrt{1-\rho^2}, &
\beta_{13}&=\alpha_{13}, &
\beta_{23}&=0 .
\end{aligned}
\]
Its eigenvalues are denoted $\theta_1\geq\theta_2\geq\theta_3$.

\paragraph{Top eigenvalue.}
Let
\[
T=\beta_{11}+\beta_{22}+\beta_{33},\qquad
Q=\beta_{11}\beta_{22}+\beta_{11}\beta_{33}+\beta_{22}\beta_{33}
-\beta_{12}^2-\beta_{13}^2,
\]
\[
D=\beta_{11}\beta_{22}\beta_{33}-\beta_{22}\beta_{13}^2-\beta_{33}\beta_{12}^2,
\qquad \Delta=T^2-3Q>0 .
\]
\[
\boxed{
\theta_1=\frac{T}{3}+\frac{2\sqrt{\Delta}}{3}
\cos\!\left(\frac13\arccos\!\left(\frac{2T^3-9TQ+27D}{2\Delta^{3/2}}\right)\right).
}
\]
Details of the computation in \ref{true}.

    
\end{comment}


\subsection{Spectrum of the Hessian $\rdmmat A$}\label{sec:hessian_A}
\subsubsection{First-step estimator}
The first-step estimator $\hat{\vect x}^{(1)}$ is the top eigenvector of $\rdmmat Y_1 = \lambda_1 \vect x_1 \vect x_1 ^T + \rdmmat Z_1$. By Theorem~\ref{thm:bbp}, its overlap with $\vect x_1$ converges to
\begin{equation}
    m_1 := \begin{cases}
        0 & \text{if } \lambda_1 \le 1,\\[0.3em]
        \sqrt{1-\lambda_1^{-2}} & \text{if } \lambda_1 > 1.
    \end{cases}
    \label{eq:lambda_max_Y1}
\end{equation}
To study the second step analytically, we replace $\hat{\vect x}^{(1)}$ by the unit vector with the same overlap,
\begin{equation}
    \hat{\vect x}^{(1)}_{P_1} = m_1\vect x_1 + \sqrt{1-m_1^2}\,\vect x_3 ,
    \label{eq:x_hat_p1}
\end{equation}
where $\vect x_3$ is a unit vector orthogonal to $\vect x_1$, independent of $\rdmmat Z_2$.

\subsubsection{The matrix}
We evaluate the Hessian at the projector estimator $\hat{\vect x}^{(1)}_{P_1}$ of Eq.~\eqref{eq:x_hat_p1}. Since $\norm{\hat{\vect x}^{(1)}_{P_1}} = 1$, the Fisher matrix of Appendix~\ref{adix:fisher} gives
\begin{equation}
    \rdmmat A = \rdmmat Y_2 - \rdmmat I_n + 2\mu_R\,\rdmmat F\big(\hat{\vect x}^{(1)}_{P_1}\big)
    = \underbrace{\Lambda\,\hat{\vect x}^{(1)}_{P_1}\hat{\vect x}^{(1)\top}_{P_1} + \lambda_2\,\vect x_2\vect x_2^\top}_{\rdmmat P}
    + \underbrace{\rdmmat Z_2}_{\rdmmat W} + c\,\rdmmat I_n,
    \qquad
    \Lambda = 2\mu_R\Big((\lambda_1-1)^2 + \frac1n\Big), \quad c = \frac{2\mu_R}{n} - 1 .
    \label{eq:A_xP1}
\end{equation}
This is the form $\rdmmat A = \rdmmat P + \rdmmat W + c\,\rdmmat I_n$ of Section~\ref{sec:model}, with $k=3$: in the orthonormal basis $\vect e_1 = \vect x_1$, $\vect e_2$, $\vect e_3 = \vect x_3$, $\rdmmat P = \rdmmat E \rdmmat B \rdmmat E^T$ with the coefficients of Table~\ref{tab:beta_A} in Appendix~\ref{tab:coefficients_beta_A(xP1)}. Its explicit form in the large $n$ limit is given in Eq.~\eqref{eq:B_A_limit}.

\subsubsection{Eigenvalues and eigenvectors}
\begin{prop}[Spectrum of $\rdmmat A$]\label{prop:spectrum_A}
Let $m_1\rho > 0$. The eigenvalues $\theta_1\ge\theta_2\ge\theta_3$ of $\rdmmat B$ are
\begin{equation}
    \theta_{1,2} = \frac{\Lambda+\lambda_2}{2} \pm \sqrt{\Big(\frac{\Lambda-\lambda_2}{2}\Big)^2 + \Lambda\lambda_2\,m_1^2\rho^2},
    \qquad \theta_3 = 0 ,
    \label{eq:theta_A}
\end{equation}
and, almost surely as $n\to\infty$, the three largest eigenvalues of $\rdmmat A$ converge to
\begin{equation}
    \lambda^{\mathrm{out}}_l \longrightarrow
    \begin{cases}
        \theta_l + \dfrac{1}{\theta_l} + c, & \theta_l > 1,\\[6pt]
        2 + c, & \theta_l \le 1,
    \end{cases}
    \quad l = 1,2,
    \qquad\qquad
    \lambda^{\mathrm{out}}_3 \longrightarrow 2 + c .
    \label{eq:outliers_A}
\end{equation}
For $l=1,2$, the eigenvectors of $\rdmmat P$ associated with $\theta_l$ and their overlaps $q_i^l = \langle\vect v_l,\vect x_i\rangle$ with the spikes are
\begin{equation}
    \vect v_l = c_1^l\,\hat{\vect x}^{(1)}_{P_1} + c_2^l\,\vect x_2, \qquad
    \big(c_1^l, c_2^l\big) = \frac{(r_l,\,1)}{\sqrt{r_l^2 + 2m_1\rho\,r_l + 1}}, \qquad
    r_l = \frac{m_1\rho\,\Lambda}{\theta_l - \Lambda} ,
    \label{eq:v_A} 
\end{equation}
\begin{equation}
    q_1^l = m_1\,c_1^l + \rho\,c_2^l, \qquad
    q_2^l = \frac{\theta_l}{\lambda_2}\,c_2^l, \qquad
    q_3^l = \sqrt{1-m_1^2}\;c_1^l ,
    \label{eq:q_A}
\end{equation}
and the unit eigenvectors $\vect u_1,\vect u_2,\vect u_3$ of $\rdmmat A$ associated with its three largest eigenvalues satisfy, for $i=1,2,3$,
\begin{equation}
    |\langle\vect u_l,\vect x_i\rangle| \longrightarrow
    \begin{cases}
        |q_i^l|\sqrt{1-\theta_l^{-2}}, & \theta_l > 1,\\
        0, & \theta_l \le 1,
    \end{cases}
    \quad l=1,2,
    \qquad
    |\langle\vect u_3,\vect x_i\rangle| \longrightarrow 0 .
    \label{eq:overlaps_A}
\end{equation}
\end{prop}
\begin{proof}
$\rdmmat P$ is built from the two vectors $\hat{\vect x}^{(1)}_{P_1}$ and $\vect x_2$, so it has rank $2$ and $\theta_3 = 0$. It has the structure of the naive signal matrix of Eq.~\eqref{eq:2spikes} with $a\to\Lambda$, $b\to\lambda_2$ and $\rho\to m_1\rho = \langle\hat{\vect x}^{(1)}_{P_1},\vect x_2\rangle$, which gives Eq.~\eqref{eq:theta_A} and, by the derivation of Appendix~\ref{sec:eigenvect_naive}, Eq.~\eqref{eq:v_A}. Eq.~\eqref{eq:q_A} follows from $\langle\vect x_1,\vect x_2\rangle = \rho$, $\langle\vect x_1,\vect x_3\rangle = 0$, $\langle\vect x_2,\vect x_3\rangle\to0$ and Eq.~\eqref{eq:overlaps-signals}. Eqs.~\eqref{eq:outliers_A} and~\eqref{eq:overlaps_A} are Theorem~\ref{thm:bbp_k} and Corollary~\ref{cor:overlap_signals}.
\end{proof}
\noindent
Since $\theta_3 = 0$, the third eigenvector belongs to the bulk and carries no information on the spikes: $\rdmmat A$ has at most two outliers. The eigenvectors and the overlaps do not depend on $c$, which only shifts the eigenvalues.

\subsubsection{Transitions}\label{sec:transitions_A}
By Eq.~\eqref{eq:outliers_A}, the number of outliers of $\rdmmat A$ is the number of $\theta_l > 1$. By Eq.~\eqref{eq:theta_A}, $\theta_1$ and $\theta_2$ are the roots of
\begin{equation}
    p(\theta) = \theta^2 - (\Lambda + \lambda_2)\,\theta + \Lambda\lambda_2(1 - m_1^2\rho^2) ,
    \label{eq:p_theta}
\end{equation}
where $\Lambda$ and $m_1 = \sqrt{1-\lambda_1^{-2}}$ (or $0$ for $\lambda_1 \le 1$) depend on $\lambda_1$ only. Solving $p(1)=0$ for $\lambda_2$ defines the two boundaries
\begin{equation}
    f_{01}(\lambda_1) = \frac{1-\Lambda}{1-\Lambda(1-m_1^2\rho^2)} \quad \text{for } \Lambda < 1,
    \qquad
    f_{12}(\lambda_1) = \frac{\Lambda-1}{\Lambda(1-m_1^2\rho^2)-1} \quad \text{for } \Lambda(1-m_1^2\rho^2) > 1 ,
    \label{eq:f01_f12}
\end{equation}
which are the same function of $\lambda_1$ on two different branches.

\begin{prop}[Phase diagram of $\rdmmat A$]\label{prop:phase_A}
In the $(\lambda_1, \lambda_2)$ plane, almost surely as $n\to\infty$, $\rdmmat A$ has
\begin{itemize}
    \item \textbf{no outlier} if $\Lambda < 1$ and $\lambda_2 < f_{01}(\lambda_1)$;
    \item \textbf{two outliers} if $\Lambda(1-m_1^2\rho^2) > 1$ and $\lambda_2 > f_{12}(\lambda_1)$;
    \item \textbf{one outlier} otherwise.
\end{itemize}
\end{prop}
\noindent
The proof is given in Appendix~\ref{adix:phase_A}. Since $f_{01} \le 1 \le f_{12}$, two outliers require $\lambda_2 > 1$, and $\lambda_2 > 1$ guarantees at least one. Equivalently, solving $p(1) = 0$ for $\Lambda$, the top eigenvalue leaves the bulk ($\theta_1 > 1$) if and only if $\lambda_2 > 1$ or $\Lambda > \Lambda^*$, with
\begin{equation}
    \Lambda^* = \frac{\lambda_2 - 1}{\lambda_2\,(1-m_1^2\rho^2) - 1},
    \qquad
    \mu_R^* = \frac{\Lambda^*}{2\big((\lambda_1-1)^2 + \frac1n\big)} .
    \label{eq:Lambda_star}
\end{equation}

\subsection{Method A: top eigenvector of the Hessian}\label{sec:methodA}
Method A takes as estimator the top eigenvector of $\rdmmat A$:
\begin{equation}
    \hat{\vect x}^{(2)}_{A} = \text{top eigenvector}(\rdmmat A) = \vect u_1 .
    \label{eq:x_hat_a}
\end{equation}
We write $\theta_{max} \equiv \theta_1$ for the largest eigenvalue of $\rdmmat B$, given by Eq.~\eqref{eq:theta_A}.

\begin{prop}[Method A]\label{prop:methodA}
Let $m_1\rho > 0$. Almost surely as $n\to\infty$:
\begin{enumerate}[(i)]
    \item the top eigenvalue of $\rdmmat A$ converges to
    \begin{equation}
        \lambda^{\mathrm{out}}_{max} \longrightarrow
        \begin{cases}
            \theta_{max} + \dfrac{1}{\theta_{max}} + \dfrac{2\mu_R}{n} - 1, & \theta_{max} > 1,\\[6pt]
            1 + \dfrac{2\mu_R}{n}, & \theta_{max} \le 1 ;
        \end{cases}
        \label{eq:top_eigenvalue_A}
    \end{equation}
    \item the estimator carries information on the signals if and only if
    \begin{equation}
        \theta_{max} > 1
        \quad\Longleftrightarrow\quad
        \lambda_2 > 1
        \quad\text{or}\quad
        \mu_R > \mu_R^* ,
        \label{eq:condition_A}
    \end{equation}
    with $\mu_R^*$ of Eq.~\eqref{eq:Lambda_star};
    \item its overlaps with the signals converge to
    \begin{equation}
        m^A_1 := |\langle \hat{\vect x}^{(2)}_A, \vect{x_1}\rangle| \longrightarrow \big\lvert m_1 c_1^1 + \rho\, c_2^1 \big\rvert \sqrt{1-\theta_{max}^{-2}} ,
        \qquad
        m^A_2 := |\langle \hat{\vect x}^{(2)}_A, \vect{x_2}\rangle| \longrightarrow \Big \lvert  \frac{\theta_{max}}{\lambda_2}c_2^1  \Big \rvert \sqrt{1-\theta_{max}^{-2}}
        \label{eq:overlap_A}
    \end{equation}
    for $\theta_{max} > 1$, and to $0$ for $\theta_{max} \le 1$, with $c_1^1, c_2^1$ of Eq.~\eqref{eq:v_A}.
\end{enumerate}
\end{prop}
\begin{proof}
(i) and (iii) are Eqs.~\eqref{eq:outliers_A} and~\eqref{eq:overlaps_A} of Proposition~\ref{prop:spectrum_A} with $l=1$, $c = \frac{2\mu_R}{n}-1$ and $q_i^1$ of Eq.~\eqref{eq:q_A}. (ii) is the condition for at least one outlier in Proposition~\ref{prop:phase_A}, written with $\Lambda^*$ of Eq.~\eqref{eq:Lambda_star}.
\end{proof}
\noindent
When the second signal is detectable on its own ($\lambda_2>1$), Method A always recovers information. Otherwise, the EWC penalty has to be strong enough to push $\theta_{max}$ above $1$: the regularisation alone can make the estimator leave the bulk. The overlaps have the same form as those of the naive method, Eq.~\eqref{eq:overlap_naive}, with $a \to \Lambda$, $b \to \lambda_2$ and $\rho \to m_1\rho$. Similarly, $m^A_3 = |\langle \hat{\vect x}^{(2)}_A, \vect{x_3}\rangle| \to \sqrt{1-m_1^2}\,|c_1^1|\sqrt{1-\theta_{max}^{-2}}$ measures the part of the first-step error inherited by the second-step estimator.

\subsection{Method B: resolvent of the Hessian}\label{sec:methodB}
Method B takes as estimator the solution of the stationarity condition with the Lagrange shift, Eq.~\eqref{eq:final_estimator_lagrange}:
\begin{equation}
    \hat{\vect x}^{(2)}_B(\lambda) = 2\mu_R\,\vect R_{\rdmmat A}(\lambda)\,\vect F_1 \hat{\vect x}^{(1)},
    \qquad \norm{\hat{\vect x}^{(2)}_B(\lambda)} = 1 ,
    \label{eq:x_hat_b}
\end{equation}
with $\hat{\vect x}^{(1)} = \hat{\vect x}^{(1)}_{P_1}$. For a unit vector $\vect x$, $\rdmmat F(\vect x)\,\vect x = \big((\lambda_1-1)^2 + \frac2n\big)\vect x$, so
\begin{equation}
    \hat{\vect x}^{(2)}_B(\lambda) = \kappa\,\vect R_{\rdmmat A}(\lambda)\,\hat{\vect x}^{(1)},
    \qquad
    \kappa = 2\mu_R\Big((\lambda_1-1)^2 + \frac2n\Big) ,
    \label{eq:xB_basisA}
\end{equation}
and the multiplier $\lambda$ is selected as the unique root of the norm constraint above the top eigenvalue $\gamma_1$ of $\rdmmat A$ (Appendix~\ref{app:B_resolvent}).

\noindent
The limit is expressed with the Stieltjes transform of the semicircle law (Appendix~\ref{proof:BBP1}), $S_W(z) = \frac{z - \sqrt{z^2-4}}{2}$ for $z>2$, which satisfies $-S_W'(z) = \frac{S_W(z)^2}{1-S_W(z)^2} > 0$, and with the overlaps $q_{il} := q_i^l$ of Eq.~\eqref{eq:q_A} and
\begin{equation}
    q_{sl} := \langle \vect v_l, \hat{\vect x}^{(1)}\rangle = \frac{\theta_l}{\Lambda}\,c^l_1 , \qquad l = 1,2,
\end{equation}
which follows from Eq.~\eqref{eq:overlaps-signals} with $a\to\Lambda$ and $\vect x_1\to\hat{\vect x}^{(1)}$. Since $\theta_3 = 0$ and $\hat{\vect x}^{(1)} \in \mathrm{span}(\vect v_1,\vect v_2)$, $q_{s3} = 0$.

\begin{prop}[Method B]\label{prop:B_overlaps}
Assume $m_1\rho > 0$ and let
$z_* := \max\!\big(2,\ \max_{l:\,\theta_l > 1}(\theta_l + \theta_l^{-1})\big)$, the right edge of the spectrum of $\rdmmat A - c\rdmmat I_n$.
Then, almost surely as $n\to\infty$, the selected multiplier converges, $\lambda \to c + z^\star$,
where $z^\star$ is the unique solution on $(z_*,\infty)$ of
\begin{equation}
    \kappa^2 \sum_{l=1}^{2} q_{sl}^2\,
    \frac{-S_W'(z)}{\big(1-\theta_l S_W(z)\big)^2} = 1 ,
    \label{eq:B_norm}
\end{equation}
and the overlaps of the estimator with the signals converge to
\begin{equation}
    m^B_i := \big|\langle \hat{\vect x}^{(2)}_B, \vect x_i\rangle\big|
    \longrightarrow
    \kappa\, \left| \sum_{l=1}^{2} q_{il}\, q_{sl}\,
    \frac{S_W(z^\star)}{1-\theta_l S_W(z^\star)} \right| ,
    \qquad i = 1,2,3 .
    \label{eq:B_overlap}
\end{equation}
\end{prop}
\noindent
The proof is given in Appendix~\ref{app:B_resolvent}. The resolvent of Eq.~\eqref{eq:xB_basisA} involves all $n$ eigenvectors of $\rdmmat A$, but its projection on the signal subspace reduces to a $2\times 2$ problem: the bulk eigenvectors are individually random, but their total contribution is captured by $S_W$. Near an outlier, $z \to \theta_l + \theta_l^{-1}$, one has $S_W(z) \to 1/\theta_l$ and
\[
    \frac{S_W(z)}{1-\theta_l S_W(z)} \simeq \frac{1-\theta_l^{-2}}{\theta_l + \theta_l^{-1} + c - \lambda} ,
\]
so the pole part of Eq.~\eqref{eq:B_overlap} is the contribution of the outlier eigenvector $\vect u_l$ given by Corollary~\ref{cor:overlap_signals}, and the regular remainder is the contribution of the bulk.

\noindent
\textbf{No sharp transition.} Unlike Method A, Method B has no recovery threshold. The left-hand side of Eq.~\eqref{eq:B_norm} diverges as $z \to z_*^+$ (through the pole $1-\theta_{\max}S_W = 0$ if $\theta_{\max} > 1$, through $-S_W' \to \infty$ as $z\to 2^+$ otherwise) and vanishes as $z\to\infty$, so a root always exists. Even for $\theta_{\max} \le 1$, Eq.~\eqref{eq:B_overlap} is non-zero as soon as $q_{il} q_{sl} \neq 0$: the estimator is then built from the bulk eigenvectors only, weighted by their overlap with $\hat{\vect x}^{(1)}$, and keeps information on $\vect x_1$, and on $\vect x_2$ through the correlation $\rho$.
